\section{Билеты}

\begin{enumerate}
\item Комплексные числа. Алгебраическая и тригонометрическая формы. Геометрический смысл операций над комплексными числами, корень.
\item Ряды с комплексными членами, свойства. Предел и непрерывность функции комплексной переменной. Сфера Римана.
\item Дифференцируемость и условия Коши—Римана.
\item Геометрический смысл модуля и аргумента производной функции \( f(z)=az,\ a\in\mathbb{C} \). Теорема о геометрическом смысле производной (постоянство растяжений и сохранение углов, \( \Gamma'(t_0)=f'(z_0)\cdot\gamma'(z_0) \)). Конформные отображения.
\item Голоморфные функции. Свойства голоморфных функций. Теорема о производной обратной функции.
\item Элементарные функции и их свойства. Геометрия экспоненты.
\item Интеграл функции комплексной переменной. Способы вычисления при заданной параметризации. Длина кривой, оценка сверху для интеграла по кривой.
\item Формула Ньютона—Лейбница. Интеграл \( \int_\gamma z^n\diff z \) при целых \( n \).
\item Теорема Коши (интеграл по границе треугольника).
\item Существование первообразной в выпуклом множестве у голоморфной функции. Теоремы Коши 2 и 3
\item Интегральная формула Коши.
\item Поточечная сходимость функциональных рядов, равномерная сходимость, свойства. Аналитичность голоморфных функций, формула Тейлора.
\item Свойства степенных рядов (Теорема Абеля, радиус сходимости, ф-ма Коши-Адамара, равномерная сходимость внутри радиуса сходимости). Теорема Тейлора.
\item Неравенство Коши. Теорема Лиувилля. Теорема о голоморфности степенного ряда
\item Свойства голоморфных функций (голоморфность производной, гармоничность Re и Im, возможность восстановить одну из частей (Re или Im) с точностью до константы функций, пример на \( u(x,y)=x^3-3xy^2 \). Все без доказательств).
\item Теорема Морера. Три разных определения голоморфных функций. Бесконечная дифференцируемость аналитических функций. Лемма о вычислении коэффициентов ряда Тейлора через производные.
\item Ряды Лорана. Теорема Лорана—Вейерштрасса
\item Изолированные особые точки и их классификация (определение всех точек и примеры). Устранимые особые точки.
\item Нули функции. Лемма о разложении голоморфной функции в нуле. Порядок нуля. Теорема о связи порядка нуля и разложении функции. Определение полюса. Теорема о характеристике полюса.
\item Классификация особых точек с помощью главной части ряда Лорана. Теорема Сохоцкого. Определение целой функции, классификация бесконечности как особой точки.
\item Определение вычета. Теорема о вычетах. Теорема о связи коэффициентов в ряде Лорана с вычетами. Вычисление вычета.
\item Вычисление вещественных интегралов при помощи вычетов.
\item Принцип аргумента, теорема Руше. Основная теорема алгебры.
\end{enumerate}